\documentclass[12pt,a4paper]{article}
\usepackage[portuguese]{babel}
\usepackage[utf8]{inputenc}
\usepackage{graphicx}
\usepackage{amsmath}
\usepackage{enumitem}
\usepackage{mathrsfs}
\usepackage{tikz}
\usepackage{gnuplot-lua-tikz}
\usepackage[framed,numbered,autolinebreaks,useliterate]{mcode}
\begin{document}
\begin{titlepage}
  \centering
  \includegraphics[width=0.33\textwidth]{usp}\par\vspace{1cm}
  {\scshape\LARGE Universidade de São Paulo\par}
  \vspace{1cm}
  {\scshape\Large SEL0611- Fundamentos de Controle\par}
  \vspace{1.5cm}
  {\huge\bfseries Lista de Exercícios No.3\par}
  \vspace{2cm}
  {\Large\itshape Pedro Morello Abbud \par}
  \vspace{1cm}
  Número USP 8058718
  \vfill
  Disciplina minsitrada por\par
  Professor Doutor B.J.Mass

  \vfill 
  % Bottom of the page
  {\large \today\par}
\end{titlepage}
\newpage
\section*{Exercícios}
Dada a equação diferencial ordinária (EDO):

\begin{align}
  y''(t)+2y'(t)+5y(t)=10x(t)
  \label{edo:1}
\end{align}
Onde $y(t)$ e $x(t)$ são respectivamente a saída e a entrada de um certo sistema, desejamos
construir uma ``imagem de raios-x” (ou ``eletrocardiograma”) do sistema, que revele a
dinâmica do mesmo. Em nosso contexto, diagramas de \emph{Bode} são verdadeiras ``imagens
de raios-x”, pois revelam coisas que não enxergamos pela mera inspeção direta da EDO.\\
Precisamos para essa visualização de algo imprescindível; um aplicativo computacional
númerico, talhado para o projeto de sistemas de controle. Duas alternativas foram apresentadas
em sala de aula: \emph{Scilab} e \emph{MATLAB}.
\begin{enumerate}
  \item Obtenha à mão, a função de transferência $G_{1}(s)= Y(s)/X(s)$ a partir 
    da EDO (\ref{edo:1}) acima.

    Sabemos que:
    \begin{align}
      x(t)=\frac{y''(t)}{10}+\frac{y'(t)}{5}+\frac{y(t)}{2} \nonumber
    \end{align} 
    Logo temos:
    \begin{align*}
      G(s)=\frac{Y(s)}{X(s)}\\
      =\frac{\mathcal{L}\{y(t)\}}{\mathcal{L}\{x(t)\}} \\
      =\frac{Y(s)}{\mathcal{L}\{\frac{y''(t)}{10}+\frac{y'(t)}{5}+\frac{y(t)}{2}\}}\\
      = \frac{10Y(s)}{Y(S)(s^2+2s+1)}\\
      = \frac{10}{s^2+2s+1}
    \end{align*}

  \item Obtenha à mão, a função de transferência $G_2(s) = X(s)/Y(s)$ a partir de (\ref{edo:1}), observando
    que, naturalmente $G_2(s) = 1/G_1(s) = G_1(s)^{-1}$

    De forma análoga ao exercício anterior sabemos que:
    \begin{align}
      x(t)=\frac{y''(t)}{10}+\frac{y'(t)}{5}+\frac{y(t)}{2} \nonumber
    \end{align} 
    Logo temos:
    \begin{align*}
      G(s)=\frac{X(s)}{Y(s)}=\frac{\mathcal{L}\{x(t)\}}{\mathcal{L}\{y(t)\}}\\
      =\frac{\mathcal{L}\{\frac{y''(t)}{10}+\frac{y'(t)}{5}+\frac{y(t)}{2}\}}{Y(s)}\\
      =  \frac{Y(S)(s^2+2s+1)}{10Y(s)}\\
      =\frac{s^2+2s+1}{10}\\
      =\frac{s^2}{10}+\frac{s}{5}+\frac{1}{10}
    \end{align*}
    Que é exatamente $G_1^{-1}$, como queríamos demonstrar.
  \item Observe que, como $G_1(s)$ é uma função de transferência \emph{``estritamente própria”},
    e  $G_2(s)$ é ``imprópria”. Os sistemas físicos de nosso interesse (engenharia de controle)
    só excepcionalmente terão funcão de transferência imprópria. Em nosso contexto 
    \emph{(Fundamentos de Controle)} consideraremos inadmissíveis funções de transferência
    impróprias, i.e.\, com numerador com grau maior que o do denominador. 
  \item Empregando o comando ``tf”, escreva uma declaração em \emph{MATLAB}, que declare
    $G_1(s)$ com o nome ``g1”. Empregue apenas minúsculas.
    \begin{lstlisting}[label=matlab:1, caption=Código em MATLAB para a função de transferência $G_1(s)$]
      g1=tf(10,[10 2 1])
    \end{lstlisting}
  \item  Consulte \emph{www.scilab.org} e verifique que comando do Scilab e análogo ou equivalente
    ao comando ``bode” do MATLAB.\\Descreva suscintamente cada uma das 6 formas
    para este comando do SCILAB.

    \begin{lstlisting}
      bode(sl, [fmin, fmax] [,step] [,comments] )
    \end{lstlisting}
    Traça um diagrama de bode tomando como argumentos de entrada: um sistema linear, os limites 
    mínimo e máximo da frequência, um degrau logarítimico descrito por um número 
    real e um vetor de strings como legenda, com a frequência dada em $hz$.
    \vfill
    \begin{lstlisting}
      bode(sl, [fmin, fmax] [,step] [,comments] [,"rad"] )
    \end{lstlisting}
    Traça um diagrama de bode tomando como argumentos de entrada: um sistema linear, os limites 
    mínimo e máximo da frequência, um degrau logarítimico descrito por um número 
    real e um vetor de strings como legenda, no entanto mostra a frequência
    em $rad/s$ ao invês de $hz$.

    \begin{lstlisting}
      bode(sl, frq [,comments] )
    \end{lstlisting}
    Traça um diagrama de bode tomando como argumentos de entrada: um sistema linear, 
    uma matriz de frequência com uma coluna para cada subsistema linear e um vetor
    de strings como legenda, com a frequência dada em $hz$

    \begin{lstlisting}
      bode(sl, frq [,comments] [,"rad"] )
    \end{lstlisting}
    Traça um diagrama de bode tomando como argumentos de entrada: um sistema linear, 
    uma matriz de frequência com uma coluna para cada subsistema linear e um vetor
    de strings como legenda, no entanto, mostra a frequência em $rad/s$ ao invês de $hz$.

    \begin{lstlisting}
      bode(frq, db, phi [,comments] )
    \end{lstlisting}
    Traça um diagram de bode tomando como argumentos de entrada: uma matriz de frequência,
    uma matriz de magnitudes dada em $Db$, uma matriz de fases dada em graus e comentários
    como legendas na forma de um vetor.

    \begin{lstlisting}
      bode(frq, repf [,comments] [,"rad"] )
    \end{lstlisting}
    Traça um diagram de bode tomando como argumentos de entrada: uma matriz de frequência,
    uma matriz de magnitudes dada em $Db$, uma matriz de fases dada em graus e comentários
    como legendas na forma de um vetor, no entanto, com a frequência em $rad/s$ ao invês de
    $hz$.
  \item Empregando a forma mais simples do comando ``bode” do MATLAB, s, imprima o
    par de gráficos ou diagramas obtidos para a funcão de transferência $G_1(s)$ acima
    referida. Retorne ao professor o par de gráficosmpressos numa única folha A4.
    \begin{lstlisting}[label=matlab:2, caption=Código em MATLAB para traçar o diagrama de Bode da função de transferência $G_1(s)$]
      g1=tf(10,[10 2 1]) 
      bode(g1)
    \end{lstlisting}
    \newpage 
    \begin{figure}
\label{fig:1}
      \begin{tikzpicture}[ scale = 1.3, transform canvas={scale=1.3}][gnuplot]
%% generated with GNUPLOT 4.6p6 (Lua 5.1; terminal rev. 99, script rev. 100)
%% 2016年03月10日 07時27分55秒
\path (0.000,0.000) rectangle (8.000,6.000);
\gpfill{rgb color={1.000,1.000,1.000}} (1.040,3.503)--(7.238,3.503)--(7.238,5.548)--(1.040,5.548)--cycle;
\gpcolor{color=gp lt color border}
\gpsetlinetype{gp lt border}
\gpsetlinewidth{1.00}
\draw[gp path] (1.040,3.503)--(1.040,5.548)--(7.238,5.548)--(7.238,3.503)--cycle;
\gpcolor{color=gp lt color axes}
\gpsetlinetype{gp lt axes}
\gpsetlinewidth{0.50}
\draw[gp path] (1.040,3.596)--(7.239,3.596);
\gpcolor{color=gp lt color border}
\gpsetlinetype{gp lt border}
\draw[gp path] (1.040,3.596)--(1.291,3.596);
\draw[gp path] (7.239,3.596)--(6.988,3.596);
\gpcolor{rgb color={0.000,0.000,0.000}}
\node[gp node right,font={\fontsize{10pt}{12pt}\selectfont}] at (0.856,3.596) {-40};
\gpcolor{color=gp lt color axes}
\gpsetlinetype{gp lt axes}
\draw[gp path] (1.040,3.884)--(7.239,3.884);
\gpcolor{color=gp lt color border}
\gpsetlinetype{gp lt border}
\draw[gp path] (1.040,3.884)--(1.291,3.884);
\draw[gp path] (7.239,3.884)--(6.988,3.884);
\gpcolor{rgb color={0.000,0.000,0.000}}
\node[gp node right,font={\fontsize{10pt}{12pt}\selectfont}] at (0.856,3.884) {-30};
\gpcolor{color=gp lt color axes}
\gpsetlinetype{gp lt axes}
\draw[gp path] (1.040,4.173)--(7.239,4.173);
\gpcolor{color=gp lt color border}
\gpsetlinetype{gp lt border}
\draw[gp path] (1.040,4.173)--(1.291,4.173);
\draw[gp path] (7.239,4.173)--(6.988,4.173);
\gpcolor{rgb color={0.000,0.000,0.000}}
\node[gp node right,font={\fontsize{10pt}{12pt}\selectfont}] at (0.856,4.173) {-20};
\gpcolor{color=gp lt color axes}
\gpsetlinetype{gp lt axes}
\draw[gp path] (1.040,4.462)--(7.239,4.462);
\gpcolor{color=gp lt color border}
\gpsetlinetype{gp lt border}
\draw[gp path] (1.040,4.462)--(1.291,4.462);
\draw[gp path] (7.239,4.462)--(6.988,4.462);
\gpcolor{rgb color={0.000,0.000,0.000}}
\node[gp node right,font={\fontsize{10pt}{12pt}\selectfont}] at (0.856,4.462) {-10};
\gpcolor{color=gp lt color axes}
\gpsetlinetype{gp lt axes}
\draw[gp path] (1.040,4.751)--(7.239,4.751);
\gpcolor{color=gp lt color border}
\gpsetlinetype{gp lt border}
\draw[gp path] (1.040,4.751)--(1.291,4.751);
\draw[gp path] (7.239,4.751)--(6.988,4.751);
\gpcolor{rgb color={0.000,0.000,0.000}}
\node[gp node right,font={\fontsize{10pt}{12pt}\selectfont}] at (0.856,4.751) {0};
\gpcolor{color=gp lt color axes}
\gpsetlinetype{gp lt axes}
\draw[gp path] (1.040,5.039)--(7.239,5.039);
\gpcolor{color=gp lt color border}
\gpsetlinetype{gp lt border}
\draw[gp path] (1.040,5.039)--(1.291,5.039);
\draw[gp path] (7.239,5.039)--(6.988,5.039);
\gpcolor{rgb color={0.000,0.000,0.000}}
\node[gp node right,font={\fontsize{10pt}{12pt}\selectfont}] at (0.856,5.039) {10};
\gpcolor{color=gp lt color axes}
\gpsetlinetype{gp lt axes}
\draw[gp path] (1.040,5.328)--(7.239,5.328);
\gpcolor{color=gp lt color border}
\gpsetlinetype{gp lt border}
\draw[gp path] (1.040,5.328)--(1.291,5.328);
\draw[gp path] (7.239,5.328)--(6.988,5.328);
\gpcolor{rgb color={0.000,0.000,0.000}}
\node[gp node right,font={\fontsize{10pt}{12pt}\selectfont}] at (0.856,5.328) {20};
\gpcolor{color=gp lt color axes}
\gpsetlinetype{gp lt axes}
\draw[gp path] (1.040,3.503)--(1.040,5.549);
\gpcolor{color=gp lt color border}
\gpsetlinetype{gp lt border}
\draw[gp path] (1.040,3.503)--(1.040,3.754);
\draw[gp path] (1.040,5.549)--(1.040,5.298);
\gpcolor{rgb color={0.000,0.000,0.000}}
\node[gp node center,font={\fontsize{10pt}{12pt}\selectfont}] at (1.040,3.195) {$10^{-2}$};
\gpcolor{color=gp lt color axes}
\gpsetlinetype{gp lt axes}
\draw[gp path] (1.662,3.503)--(1.662,5.549);
\gpcolor{color=gp lt color border}
\gpsetlinetype{gp lt border}
\draw[gp path] (1.662,3.503)--(1.662,3.628);
\draw[gp path] (1.662,5.549)--(1.662,5.424);
\gpcolor{color=gp lt color axes}
\gpsetlinetype{gp lt axes}
\draw[gp path] (2.026,3.503)--(2.026,5.549);
\gpcolor{color=gp lt color border}
\gpsetlinetype{gp lt border}
\draw[gp path] (2.026,3.503)--(2.026,3.628);
\draw[gp path] (2.026,5.549)--(2.026,5.424);
\gpcolor{color=gp lt color axes}
\gpsetlinetype{gp lt axes}
\draw[gp path] (2.284,3.503)--(2.284,5.549);
\gpcolor{color=gp lt color border}
\gpsetlinetype{gp lt border}
\draw[gp path] (2.284,3.503)--(2.284,3.628);
\draw[gp path] (2.284,5.549)--(2.284,5.424);
\gpcolor{color=gp lt color axes}
\gpsetlinetype{gp lt axes}
\draw[gp path] (2.484,3.503)--(2.484,5.549);
\gpcolor{color=gp lt color border}
\gpsetlinetype{gp lt border}
\draw[gp path] (2.484,3.503)--(2.484,3.628);
\draw[gp path] (2.484,5.549)--(2.484,5.424);
\gpcolor{color=gp lt color axes}
\gpsetlinetype{gp lt axes}
\draw[gp path] (2.648,3.503)--(2.648,5.549);
\gpcolor{color=gp lt color border}
\gpsetlinetype{gp lt border}
\draw[gp path] (2.648,3.503)--(2.648,3.628);
\draw[gp path] (2.648,5.549)--(2.648,5.424);
\gpcolor{color=gp lt color axes}
\gpsetlinetype{gp lt axes}
\draw[gp path] (2.786,3.503)--(2.786,5.549);
\gpcolor{color=gp lt color border}
\gpsetlinetype{gp lt border}
\draw[gp path] (2.786,3.503)--(2.786,3.628);
\draw[gp path] (2.786,5.549)--(2.786,5.424);
\gpcolor{color=gp lt color axes}
\gpsetlinetype{gp lt axes}
\draw[gp path] (2.906,3.503)--(2.906,5.549);
\gpcolor{color=gp lt color border}
\gpsetlinetype{gp lt border}
\draw[gp path] (2.906,3.503)--(2.906,3.628);
\draw[gp path] (2.906,5.549)--(2.906,5.424);
\gpcolor{color=gp lt color axes}
\gpsetlinetype{gp lt axes}
\draw[gp path] (3.012,3.503)--(3.012,5.549);
\gpcolor{color=gp lt color border}
\gpsetlinetype{gp lt border}
\draw[gp path] (3.012,3.503)--(3.012,3.628);
\draw[gp path] (3.012,5.549)--(3.012,5.424);
\gpcolor{color=gp lt color axes}
\gpsetlinetype{gp lt axes}
\draw[gp path] (3.106,3.503)--(3.106,5.549);
\gpcolor{color=gp lt color border}
\gpsetlinetype{gp lt border}
\draw[gp path] (3.106,3.503)--(3.106,3.754);
\draw[gp path] (3.106,5.549)--(3.106,5.298);
\gpcolor{rgb color={0.000,0.000,0.000}}
\node[gp node center,font={\fontsize{10pt}{12pt}\selectfont}] at (3.106,3.195) {$10^{-1}$};
\gpcolor{color=gp lt color axes}
\gpsetlinetype{gp lt axes}
\draw[gp path] (3.728,3.503)--(3.728,5.549);
\gpcolor{color=gp lt color border}
\gpsetlinetype{gp lt border}
\draw[gp path] (3.728,3.503)--(3.728,3.628);
\draw[gp path] (3.728,5.549)--(3.728,5.424);
\gpcolor{color=gp lt color axes}
\gpsetlinetype{gp lt axes}
\draw[gp path] (4.092,3.503)--(4.092,5.549);
\gpcolor{color=gp lt color border}
\gpsetlinetype{gp lt border}
\draw[gp path] (4.092,3.503)--(4.092,3.628);
\draw[gp path] (4.092,5.549)--(4.092,5.424);
\gpcolor{color=gp lt color axes}
\gpsetlinetype{gp lt axes}
\draw[gp path] (4.350,3.503)--(4.350,5.549);
\gpcolor{color=gp lt color border}
\gpsetlinetype{gp lt border}
\draw[gp path] (4.350,3.503)--(4.350,3.628);
\draw[gp path] (4.350,5.549)--(4.350,5.424);
\gpcolor{color=gp lt color axes}
\gpsetlinetype{gp lt axes}
\draw[gp path] (4.551,3.503)--(4.551,5.549);
\gpcolor{color=gp lt color border}
\gpsetlinetype{gp lt border}
\draw[gp path] (4.551,3.503)--(4.551,3.628);
\draw[gp path] (4.551,5.549)--(4.551,5.424);
\gpcolor{color=gp lt color axes}
\gpsetlinetype{gp lt axes}
\draw[gp path] (4.714,3.503)--(4.714,5.549);
\gpcolor{color=gp lt color border}
\gpsetlinetype{gp lt border}
\draw[gp path] (4.714,3.503)--(4.714,3.628);
\draw[gp path] (4.714,5.549)--(4.714,5.424);
\gpcolor{color=gp lt color axes}
\gpsetlinetype{gp lt axes}
\draw[gp path] (4.853,3.503)--(4.853,5.549);
\gpcolor{color=gp lt color border}
\gpsetlinetype{gp lt border}
\draw[gp path] (4.853,3.503)--(4.853,3.628);
\draw[gp path] (4.853,5.549)--(4.853,5.424);
\gpcolor{color=gp lt color axes}
\gpsetlinetype{gp lt axes}
\draw[gp path] (4.972,3.503)--(4.972,5.549);
\gpcolor{color=gp lt color border}
\gpsetlinetype{gp lt border}
\draw[gp path] (4.972,3.503)--(4.972,3.628);
\draw[gp path] (4.972,5.549)--(4.972,5.424);
\gpcolor{color=gp lt color axes}
\gpsetlinetype{gp lt axes}
\draw[gp path] (5.078,3.503)--(5.078,5.549);
\gpcolor{color=gp lt color border}
\gpsetlinetype{gp lt border}
\draw[gp path] (5.078,3.503)--(5.078,3.628);
\draw[gp path] (5.078,5.549)--(5.078,5.424);
\gpcolor{color=gp lt color axes}
\gpsetlinetype{gp lt axes}
\draw[gp path] (5.173,3.503)--(5.173,5.549);
\gpcolor{color=gp lt color border}
\gpsetlinetype{gp lt border}
\draw[gp path] (5.173,3.503)--(5.173,3.754);
\draw[gp path] (5.173,5.549)--(5.173,5.298);
\gpcolor{rgb color={0.000,0.000,0.000}}
\node[gp node center,font={\fontsize{10pt}{12pt}\selectfont}] at (5.173,3.195) {$10^{0}$};
\gpcolor{color=gp lt color axes}
\gpsetlinetype{gp lt axes}
\draw[gp path] (5.795,3.503)--(5.795,5.549);
\gpcolor{color=gp lt color border}
\gpsetlinetype{gp lt border}
\draw[gp path] (5.795,3.503)--(5.795,3.628);
\draw[gp path] (5.795,5.549)--(5.795,5.424);
\gpcolor{color=gp lt color axes}
\gpsetlinetype{gp lt axes}
\draw[gp path] (6.159,3.503)--(6.159,5.549);
\gpcolor{color=gp lt color border}
\gpsetlinetype{gp lt border}
\draw[gp path] (6.159,3.503)--(6.159,3.628);
\draw[gp path] (6.159,5.549)--(6.159,5.424);
\gpcolor{color=gp lt color axes}
\gpsetlinetype{gp lt axes}
\draw[gp path] (6.417,3.503)--(6.417,5.549);
\gpcolor{color=gp lt color border}
\gpsetlinetype{gp lt border}
\draw[gp path] (6.417,3.503)--(6.417,3.628);
\draw[gp path] (6.417,5.549)--(6.417,5.424);
\gpcolor{color=gp lt color axes}
\gpsetlinetype{gp lt axes}
\draw[gp path] (6.617,3.503)--(6.617,5.549);
\gpcolor{color=gp lt color border}
\gpsetlinetype{gp lt border}
\draw[gp path] (6.617,3.503)--(6.617,3.628);
\draw[gp path] (6.617,5.549)--(6.617,5.424);
\gpcolor{color=gp lt color axes}
\gpsetlinetype{gp lt axes}
\draw[gp path] (6.781,3.503)--(6.781,5.549);
\gpcolor{color=gp lt color border}
\gpsetlinetype{gp lt border}
\draw[gp path] (6.781,3.503)--(6.781,3.628);
\draw[gp path] (6.781,5.549)--(6.781,5.424);
\gpcolor{color=gp lt color axes}
\gpsetlinetype{gp lt axes}
\draw[gp path] (6.919,3.503)--(6.919,5.549);
\gpcolor{color=gp lt color border}
\gpsetlinetype{gp lt border}
\draw[gp path] (6.919,3.503)--(6.919,3.628);
\draw[gp path] (6.919,5.549)--(6.919,5.424);
\gpcolor{color=gp lt color axes}
\gpsetlinetype{gp lt axes}
\draw[gp path] (7.039,3.503)--(7.039,5.549);
\gpcolor{color=gp lt color border}
\gpsetlinetype{gp lt border}
\draw[gp path] (7.039,3.503)--(7.039,3.628);
\draw[gp path] (7.039,5.549)--(7.039,5.424);
\gpcolor{color=gp lt color axes}
\gpsetlinetype{gp lt axes}
\draw[gp path] (7.144,3.503)--(7.144,5.549);
\gpcolor{color=gp lt color border}
\gpsetlinetype{gp lt border}
\draw[gp path] (7.144,3.503)--(7.144,3.628);
\draw[gp path] (7.144,5.549)--(7.144,5.424);
\gpcolor{color=gp lt color axes}
\gpsetlinetype{gp lt axes}
\draw[gp path] (7.239,3.503)--(7.239,5.549);
\gpcolor{color=gp lt color border}
\gpsetlinetype{gp lt border}
\draw[gp path] (7.239,3.503)--(7.239,3.754);
\draw[gp path] (7.239,5.549)--(7.239,5.298);
\gpcolor{rgb color={0.000,0.000,0.000}}
\node[gp node center,font={\fontsize{10pt}{12pt}\selectfont}] at (7.239,3.195) {$10^{1}$};
\gpcolor{color=gp lt color border}
\draw[gp path] (1.040,5.549)--(1.040,3.503)--(7.239,3.503)--(7.239,5.549)--cycle;
\gpcolor{rgb color={0.000,0.000,0.000}}
\node[gp node center,rotate=90,font={\fontsize{10pt}{12pt}\selectfont}] at (-0.034,4.526) {Magnitude [dB]};
\gpcolor{color=gp lt color border}
\node[gp node center,font={\fontsize{10pt}{12pt}\selectfont}] at (4.139,6.011) {Bode Diagram};
\gpcolor{rgb color={0.000,0.000,1.000}}
\gpsetlinetype{gp lt plot 0}
\draw[gp path] (1.040,5.328)--(1.052,5.328)--(1.065,5.328)--(1.077,5.328)--(1.090,5.328)%
  --(1.102,5.328)--(1.115,5.328)--(1.127,5.328)--(1.139,5.328)--(1.152,5.328)--(1.164,5.328)%
  --(1.177,5.328)--(1.189,5.328)--(1.201,5.328)--(1.214,5.328)--(1.226,5.328)--(1.239,5.328)%
  --(1.251,5.328)--(1.264,5.328)--(1.276,5.328)--(1.288,5.328)--(1.301,5.328)--(1.313,5.328)%
  --(1.326,5.328)--(1.338,5.328)--(1.351,5.328)--(1.363,5.328)--(1.375,5.328)--(1.388,5.328)%
  --(1.400,5.328)--(1.413,5.328)--(1.425,5.328)--(1.438,5.328)--(1.450,5.328)--(1.462,5.328)%
  --(1.475,5.328)--(1.487,5.328)--(1.500,5.329)--(1.512,5.329)--(1.524,5.329)--(1.537,5.329)%
  --(1.549,5.329)--(1.562,5.329)--(1.574,5.329)--(1.587,5.329)--(1.599,5.329)--(1.611,5.329)%
  --(1.624,5.329)--(1.636,5.329)--(1.649,5.329)--(1.661,5.329)--(1.674,5.329)--(1.686,5.329)%
  --(1.698,5.329)--(1.711,5.329)--(1.723,5.329)--(1.736,5.329)--(1.748,5.329)--(1.761,5.329)%
  --(1.773,5.329)--(1.785,5.329)--(1.798,5.329)--(1.810,5.329)--(1.823,5.329)--(1.835,5.329)%
  --(1.847,5.329)--(1.860,5.329)--(1.872,5.329)--(1.885,5.329)--(1.897,5.329)--(1.910,5.329)%
  --(1.922,5.329)--(1.934,5.329)--(1.947,5.329)--(1.959,5.330)--(1.972,5.330)--(1.984,5.330)%
  --(1.997,5.330)--(2.009,5.330)--(2.021,5.330)--(2.034,5.330)--(2.046,5.330)--(2.059,5.330)%
  --(2.071,5.330)--(2.084,5.330)--(2.096,5.330)--(2.108,5.330)--(2.121,5.330)--(2.133,5.330)%
  --(2.146,5.330)--(2.158,5.330)--(2.170,5.330)--(2.183,5.331)--(2.195,5.331)--(2.208,5.331)%
  --(2.220,5.331)--(2.233,5.331)--(2.245,5.331)--(2.257,5.331)--(2.270,5.331)--(2.282,5.331)%
  --(2.295,5.331)--(2.307,5.331)--(2.320,5.331)--(2.332,5.332)--(2.344,5.332)--(2.357,5.332)%
  --(2.369,5.332)--(2.382,5.332)--(2.394,5.332)--(2.407,5.332)--(2.419,5.332)--(2.431,5.332)%
  --(2.444,5.333)--(2.456,5.333)--(2.469,5.333)--(2.481,5.333)--(2.493,5.333)--(2.506,5.333)%
  --(2.518,5.333)--(2.531,5.334)--(2.543,5.334)--(2.556,5.334)--(2.568,5.334)--(2.580,5.334)%
  --(2.593,5.334)--(2.605,5.335)--(2.618,5.335)--(2.630,5.335)--(2.643,5.335)--(2.655,5.335)%
  --(2.667,5.336)--(2.680,5.336)--(2.692,5.336)--(2.705,5.336)--(2.717,5.336)--(2.730,5.337)%
  --(2.742,5.337)--(2.754,5.337)--(2.767,5.337)--(2.779,5.338)--(2.792,5.338)--(2.804,5.338)%
  --(2.816,5.339)--(2.829,5.339)--(2.841,5.339)--(2.854,5.339)--(2.866,5.340)--(2.879,5.340)%
  --(2.891,5.340)--(2.903,5.341)--(2.916,5.341)--(2.928,5.342)--(2.941,5.342)--(2.953,5.342)%
  --(2.966,5.343)--(2.978,5.343)--(2.990,5.344)--(3.003,5.344)--(3.015,5.345)--(3.028,5.345)%
  --(3.040,5.345)--(3.053,5.346)--(3.065,5.346)--(3.077,5.347)--(3.090,5.348)--(3.102,5.348)%
  --(3.115,5.349)--(3.127,5.349)--(3.139,5.350)--(3.152,5.351)--(3.164,5.351)--(3.177,5.352)%
  --(3.189,5.353)--(3.202,5.353)--(3.214,5.354)--(3.226,5.355)--(3.239,5.355)--(3.251,5.356)%
  --(3.264,5.357)--(3.276,5.358)--(3.289,5.359)--(3.301,5.360)--(3.313,5.360)--(3.326,5.361)%
  --(3.338,5.362)--(3.351,5.363)--(3.363,5.364)--(3.375,5.365)--(3.388,5.366)--(3.400,5.368)%
  --(3.413,5.369)--(3.425,5.370)--(3.438,5.371)--(3.450,5.372)--(3.462,5.374)--(3.475,5.375)%
  --(3.487,5.376)--(3.500,5.378)--(3.512,5.379)--(3.525,5.380)--(3.537,5.382)--(3.549,5.383)%
  --(3.562,5.385)--(3.574,5.387)--(3.587,5.388)--(3.599,5.390)--(3.612,5.392)--(3.624,5.393)%
  --(3.636,5.395)--(3.649,5.397)--(3.661,5.399)--(3.674,5.401)--(3.686,5.403)--(3.698,5.405)%
  --(3.711,5.407)--(3.723,5.409)--(3.736,5.411)--(3.748,5.413)--(3.761,5.416)--(3.773,5.418)%
  --(3.785,5.420)--(3.798,5.422)--(3.810,5.425)--(3.823,5.427)--(3.835,5.429)--(3.848,5.432)%
  --(3.860,5.434)--(3.872,5.436)--(3.885,5.438)--(3.897,5.441)--(3.910,5.443)--(3.922,5.445)%
  --(3.935,5.447)--(3.947,5.449)--(3.959,5.450)--(3.972,5.452)--(3.984,5.453)--(3.997,5.454)%
  --(4.009,5.455)--(4.021,5.456)--(4.034,5.456)--(4.046,5.456)--(4.059,5.456)--(4.071,5.455)%
  --(4.084,5.454)--(4.096,5.452)--(4.108,5.450)--(4.121,5.447)--(4.133,5.444)--(4.139,5.443)%
  --(4.146,5.441)--(4.158,5.437)--(4.171,5.433)--(4.183,5.428)--(4.195,5.422)--(4.208,5.417)%
  --(4.220,5.410)--(4.233,5.404)--(4.245,5.397)--(4.258,5.390)--(4.270,5.382)--(4.282,5.374)%
  --(4.295,5.366)--(4.307,5.358)--(4.320,5.350)--(4.332,5.341)--(4.344,5.332)--(4.357,5.323)%
  --(4.369,5.314)--(4.382,5.305)--(4.394,5.296)--(4.407,5.287)--(4.419,5.278)--(4.431,5.268)%
  --(4.444,5.259)--(4.456,5.250)--(4.469,5.241)--(4.481,5.231)--(4.494,5.222)--(4.506,5.213)%
  --(4.518,5.204)--(4.531,5.195)--(4.543,5.186)--(4.556,5.176)--(4.568,5.167)--(4.581,5.158)%
  --(4.593,5.149)--(4.605,5.141)--(4.618,5.132)--(4.630,5.123)--(4.643,5.114)--(4.655,5.105)%
  --(4.667,5.097)--(4.680,5.088)--(4.692,5.079)--(4.705,5.071)--(4.717,5.062)--(4.730,5.054)%
  --(4.742,5.045)--(4.754,5.037)--(4.767,5.028)--(4.779,5.020)--(4.792,5.012)--(4.804,5.003)%
  --(4.817,4.995)--(4.829,4.987)--(4.841,4.979)--(4.854,4.971)--(4.866,4.963)--(4.879,4.954)%
  --(4.891,4.946)--(4.904,4.938)--(4.916,4.930)--(4.928,4.923)--(4.941,4.915)--(4.953,4.907)%
  --(4.966,4.899)--(4.978,4.891)--(4.990,4.883)--(5.003,4.875)--(5.015,4.868)--(5.028,4.860)%
  --(5.040,4.852)--(5.053,4.844)--(5.065,4.837)--(5.077,4.829)--(5.090,4.821)--(5.102,4.814)%
  --(5.115,4.806)--(5.127,4.799)--(5.140,4.791)--(5.152,4.783)--(5.164,4.776)--(5.177,4.768)%
  --(5.189,4.761)--(5.202,4.753)--(5.214,4.746)--(5.226,4.739)--(5.239,4.731)--(5.251,4.724)%
  --(5.264,4.716)--(5.276,4.709)--(5.289,4.701)--(5.301,4.694)--(5.313,4.687)--(5.326,4.679)%
  --(5.338,4.672)--(5.351,4.665)--(5.363,4.657)--(5.376,4.650)--(5.388,4.643)--(5.400,4.635)%
  --(5.413,4.628)--(5.425,4.621)--(5.438,4.614)--(5.450,4.606)--(5.463,4.599)--(5.475,4.592)%
  --(5.487,4.585)--(5.500,4.577)--(5.512,4.570)--(5.525,4.563)--(5.537,4.556)--(5.549,4.549)%
  --(5.562,4.542)--(5.574,4.534)--(5.587,4.527)--(5.599,4.520)--(5.612,4.513)--(5.624,4.506)%
  --(5.636,4.499)--(5.649,4.491)--(5.661,4.484)--(5.674,4.477)--(5.686,4.470)--(5.699,4.463)%
  --(5.711,4.456)--(5.723,4.449)--(5.736,4.442)--(5.748,4.434)--(5.761,4.427)--(5.773,4.420)%
  --(5.786,4.413)--(5.798,4.406)--(5.810,4.399)--(5.823,4.392)--(5.835,4.385)--(5.848,4.378)%
  --(5.860,4.371)--(5.872,4.364)--(5.885,4.357)--(5.897,4.350)--(5.910,4.343)--(5.922,4.335)%
  --(5.935,4.328)--(5.947,4.321)--(5.959,4.314)--(5.972,4.307)--(5.984,4.300)--(5.997,4.293)%
  --(6.009,4.286)--(6.022,4.279)--(6.034,4.272)--(6.046,4.265)--(6.059,4.258)--(6.071,4.251)%
  --(6.084,4.244)--(6.096,4.237)--(6.109,4.230)--(6.121,4.223)--(6.133,4.216)--(6.146,4.209)%
  --(6.158,4.202)--(6.171,4.195)--(6.183,4.188)--(6.195,4.181)--(6.208,4.174)--(6.220,4.167)%
  --(6.233,4.160)--(6.245,4.153)--(6.258,4.146)--(6.270,4.139)--(6.282,4.132)--(6.295,4.125)%
  --(6.307,4.118)--(6.320,4.111)--(6.332,4.104)--(6.345,4.097)--(6.357,4.090)--(6.369,4.083)%
  --(6.382,4.076)--(6.394,4.069)--(6.407,4.062)--(6.419,4.055)--(6.432,4.048)--(6.444,4.041)%
  --(6.456,4.034)--(6.469,4.027)--(6.481,4.020)--(6.494,4.013)--(6.506,4.006)--(6.518,3.999)%
  --(6.531,3.992)--(6.543,3.986)--(6.556,3.979)--(6.568,3.972)--(6.581,3.965)--(6.593,3.958)%
  --(6.605,3.951)--(6.618,3.944)--(6.630,3.937)--(6.643,3.930)--(6.655,3.923)--(6.668,3.916)%
  --(6.680,3.909)--(6.692,3.902)--(6.705,3.895)--(6.717,3.888)--(6.730,3.881)--(6.742,3.874)%
  --(6.755,3.867)--(6.767,3.860)--(6.779,3.853)--(6.792,3.846)--(6.804,3.839)--(6.817,3.832)%
  --(6.829,3.825)--(6.841,3.818)--(6.854,3.811)--(6.866,3.805)--(6.879,3.798)--(6.891,3.791)%
  --(6.904,3.784)--(6.916,3.777)--(6.928,3.770)--(6.941,3.763)--(6.953,3.756)--(6.966,3.749)%
  --(6.978,3.742)--(6.991,3.735)--(7.003,3.728)--(7.015,3.721)--(7.028,3.714)--(7.040,3.707)%
  --(7.053,3.700)--(7.065,3.693)--(7.078,3.686)--(7.090,3.679)--(7.102,3.672)--(7.115,3.665)%
  --(7.127,3.659)--(7.140,3.652)--(7.152,3.645)--(7.164,3.638)--(7.177,3.631)--(7.189,3.624)%
  --(7.202,3.617)--(7.214,3.610)--(7.227,3.603)--(7.239,3.596);
%% coordinates of the plot area
\gpdefrectangularnode{gp plot 1}{\pgfpoint{1.040cm}{3.503cm}}{\pgfpoint{7.239cm}{5.549cm}}
\gpfill{rgb color={1.000,1.000,1.000}} (1.040,0.660)--(7.238,0.660)--(7.238,2.705)--(1.040,2.705)--cycle;
\gpcolor{color=gp lt color border}
\gpsetlinetype{gp lt border}
\gpsetlinewidth{1.00}
\draw[gp path] (1.040,0.660)--(1.040,2.705)--(7.238,2.705)--(7.238,0.660)--cycle;
\gpcolor{color=gp lt color axes}
\gpsetlinetype{gp lt axes}
\gpsetlinewidth{0.50}
\draw[gp path] (1.040,0.741)--(7.239,0.741);
\gpcolor{color=gp lt color border}
\gpsetlinetype{gp lt border}
\draw[gp path] (1.040,0.741)--(1.291,0.741);
\draw[gp path] (7.239,0.741)--(6.988,0.741);
\gpcolor{rgb color={0.000,0.000,0.000}}
\node[gp node right,font={\fontsize{10pt}{12pt}\selectfont}] at (0.856,0.741) {-180};
\gpcolor{color=gp lt color axes}
\gpsetlinetype{gp lt axes}
\draw[gp path] (1.040,0.950)--(7.239,0.950);
\gpcolor{color=gp lt color border}
\gpsetlinetype{gp lt border}
\draw[gp path] (1.040,0.950)--(1.291,0.950);
\draw[gp path] (7.239,0.950)--(6.988,0.950);
\gpcolor{rgb color={0.000,0.000,0.000}}
\node[gp node right,font={\fontsize{10pt}{12pt}\selectfont}] at (0.856,0.950) {-160};
\gpcolor{color=gp lt color axes}
\gpsetlinetype{gp lt axes}
\draw[gp path] (1.040,1.160)--(7.239,1.160);
\gpcolor{color=gp lt color border}
\gpsetlinetype{gp lt border}
\draw[gp path] (1.040,1.160)--(1.291,1.160);
\draw[gp path] (7.239,1.160)--(6.988,1.160);
\gpcolor{rgb color={0.000,0.000,0.000}}
\node[gp node right,font={\fontsize{10pt}{12pt}\selectfont}] at (0.856,1.160) {-140};
\gpcolor{color=gp lt color axes}
\gpsetlinetype{gp lt axes}
\draw[gp path] (1.040,1.369)--(7.239,1.369);
\gpcolor{color=gp lt color border}
\gpsetlinetype{gp lt border}
\draw[gp path] (1.040,1.369)--(1.291,1.369);
\draw[gp path] (7.239,1.369)--(6.988,1.369);
\gpcolor{rgb color={0.000,0.000,0.000}}
\node[gp node right,font={\fontsize{10pt}{12pt}\selectfont}] at (0.856,1.369) {-120};
\gpcolor{color=gp lt color axes}
\gpsetlinetype{gp lt axes}
\draw[gp path] (1.040,1.578)--(7.239,1.578);
\gpcolor{color=gp lt color border}
\gpsetlinetype{gp lt border}
\draw[gp path] (1.040,1.578)--(1.291,1.578);
\draw[gp path] (7.239,1.578)--(6.988,1.578);
\gpcolor{rgb color={0.000,0.000,0.000}}
\node[gp node right,font={\fontsize{10pt}{12pt}\selectfont}] at (0.856,1.578) {-100};
\gpcolor{color=gp lt color axes}
\gpsetlinetype{gp lt axes}
\draw[gp path] (1.040,1.788)--(7.239,1.788);
\gpcolor{color=gp lt color border}
\gpsetlinetype{gp lt border}
\draw[gp path] (1.040,1.788)--(1.291,1.788);
\draw[gp path] (7.239,1.788)--(6.988,1.788);
\gpcolor{rgb color={0.000,0.000,0.000}}
\node[gp node right,font={\fontsize{10pt}{12pt}\selectfont}] at (0.856,1.788) {-80};
\gpcolor{color=gp lt color axes}
\gpsetlinetype{gp lt axes}
\draw[gp path] (1.040,1.997)--(7.239,1.997);
\gpcolor{color=gp lt color border}
\gpsetlinetype{gp lt border}
\draw[gp path] (1.040,1.997)--(1.291,1.997);
\draw[gp path] (7.239,1.997)--(6.988,1.997);
\gpcolor{rgb color={0.000,0.000,0.000}}
\node[gp node right,font={\fontsize{10pt}{12pt}\selectfont}] at (0.856,1.997) {-60};
\gpcolor{color=gp lt color axes}
\gpsetlinetype{gp lt axes}
\draw[gp path] (1.040,2.206)--(7.239,2.206);
\gpcolor{color=gp lt color border}
\gpsetlinetype{gp lt border}
\draw[gp path] (1.040,2.206)--(1.291,2.206);
\draw[gp path] (7.239,2.206)--(6.988,2.206);
\gpcolor{rgb color={0.000,0.000,0.000}}
\node[gp node right,font={\fontsize{10pt}{12pt}\selectfont}] at (0.856,2.206) {-40};
\gpcolor{color=gp lt color axes}
\gpsetlinetype{gp lt axes}
\draw[gp path] (1.040,2.416)--(5.403,2.416);
\draw[gp path] (7.055,2.416)--(7.239,2.416);
\gpcolor{color=gp lt color border}
\gpsetlinetype{gp lt border}
\draw[gp path] (1.040,2.416)--(1.291,2.416);
\draw[gp path] (7.239,2.416)--(6.988,2.416);
\gpcolor{rgb color={0.000,0.000,0.000}}
\node[gp node right,font={\fontsize{10pt}{12pt}\selectfont}] at (0.856,2.416) {-20};
\gpcolor{color=gp lt color axes}
\gpsetlinetype{gp lt axes}
\draw[gp path] (1.040,2.625)--(7.239,2.625);
\gpcolor{color=gp lt color border}
\gpsetlinetype{gp lt border}
\draw[gp path] (1.040,2.625)--(1.291,2.625);
\draw[gp path] (7.239,2.625)--(6.988,2.625);
\gpcolor{rgb color={0.000,0.000,0.000}}
\node[gp node right,font={\fontsize{10pt}{12pt}\selectfont}] at (0.856,2.625) {0};
\gpcolor{color=gp lt color axes}
\gpsetlinetype{gp lt axes}
\draw[gp path] (1.040,0.660)--(1.040,2.706);
\gpcolor{color=gp lt color border}
\gpsetlinetype{gp lt border}
\draw[gp path] (1.040,0.660)--(1.040,0.911);
\draw[gp path] (1.040,2.706)--(1.040,2.455);
\gpcolor{rgb color={0.000,0.000,0.000}}
\node[gp node center,font={\fontsize{10pt}{12pt}\selectfont}] at (1.040,0.352) {$10^{-2}$};
\gpcolor{color=gp lt color axes}
\gpsetlinetype{gp lt axes}
\draw[gp path] (1.662,0.660)--(1.662,2.706);
\gpcolor{color=gp lt color border}
\gpsetlinetype{gp lt border}
\draw[gp path] (1.662,0.660)--(1.662,0.785);
\draw[gp path] (1.662,2.706)--(1.662,2.581);
\gpcolor{color=gp lt color axes}
\gpsetlinetype{gp lt axes}
\draw[gp path] (2.026,0.660)--(2.026,2.706);
\gpcolor{color=gp lt color border}
\gpsetlinetype{gp lt border}
\draw[gp path] (2.026,0.660)--(2.026,0.785);
\draw[gp path] (2.026,2.706)--(2.026,2.581);
\gpcolor{color=gp lt color axes}
\gpsetlinetype{gp lt axes}
\draw[gp path] (2.284,0.660)--(2.284,2.706);
\gpcolor{color=gp lt color border}
\gpsetlinetype{gp lt border}
\draw[gp path] (2.284,0.660)--(2.284,0.785);
\draw[gp path] (2.284,2.706)--(2.284,2.581);
\gpcolor{color=gp lt color axes}
\gpsetlinetype{gp lt axes}
\draw[gp path] (2.484,0.660)--(2.484,2.706);
\gpcolor{color=gp lt color border}
\gpsetlinetype{gp lt border}
\draw[gp path] (2.484,0.660)--(2.484,0.785);
\draw[gp path] (2.484,2.706)--(2.484,2.581);
\gpcolor{color=gp lt color axes}
\gpsetlinetype{gp lt axes}
\draw[gp path] (2.648,0.660)--(2.648,2.706);
\gpcolor{color=gp lt color border}
\gpsetlinetype{gp lt border}
\draw[gp path] (2.648,0.660)--(2.648,0.785);
\draw[gp path] (2.648,2.706)--(2.648,2.581);
\gpcolor{color=gp lt color axes}
\gpsetlinetype{gp lt axes}
\draw[gp path] (2.786,0.660)--(2.786,2.706);
\gpcolor{color=gp lt color border}
\gpsetlinetype{gp lt border}
\draw[gp path] (2.786,0.660)--(2.786,0.785);
\draw[gp path] (2.786,2.706)--(2.786,2.581);
\gpcolor{color=gp lt color axes}
\gpsetlinetype{gp lt axes}
\draw[gp path] (2.906,0.660)--(2.906,2.706);
\gpcolor{color=gp lt color border}
\gpsetlinetype{gp lt border}
\draw[gp path] (2.906,0.660)--(2.906,0.785);
\draw[gp path] (2.906,2.706)--(2.906,2.581);
\gpcolor{color=gp lt color axes}
\gpsetlinetype{gp lt axes}
\draw[gp path] (3.012,0.660)--(3.012,2.706);
\gpcolor{color=gp lt color border}
\gpsetlinetype{gp lt border}
\draw[gp path] (3.012,0.660)--(3.012,0.785);
\draw[gp path] (3.012,2.706)--(3.012,2.581);
\gpcolor{color=gp lt color axes}
\gpsetlinetype{gp lt axes}
\draw[gp path] (3.106,0.660)--(3.106,2.706);
\gpcolor{color=gp lt color border}
\gpsetlinetype{gp lt border}
\draw[gp path] (3.106,0.660)--(3.106,0.911);
\draw[gp path] (3.106,2.706)--(3.106,2.455);
\gpcolor{rgb color={0.000,0.000,0.000}}
\node[gp node center,font={\fontsize{10pt}{12pt}\selectfont}] at (3.106,0.352) {$10^{-1}$};
\gpcolor{color=gp lt color axes}
\gpsetlinetype{gp lt axes}
\draw[gp path] (3.728,0.660)--(3.728,2.706);
\gpcolor{color=gp lt color border}
\gpsetlinetype{gp lt border}
\draw[gp path] (3.728,0.660)--(3.728,0.785);
\draw[gp path] (3.728,2.706)--(3.728,2.581);
\gpcolor{color=gp lt color axes}
\gpsetlinetype{gp lt axes}
\draw[gp path] (4.092,0.660)--(4.092,2.706);
\gpcolor{color=gp lt color border}
\gpsetlinetype{gp lt border}
\draw[gp path] (4.092,0.660)--(4.092,0.785);
\draw[gp path] (4.092,2.706)--(4.092,2.581);
\gpcolor{color=gp lt color axes}
\gpsetlinetype{gp lt axes}
\draw[gp path] (4.350,0.660)--(4.350,2.706);
\gpcolor{color=gp lt color border}
\gpsetlinetype{gp lt border}
\draw[gp path] (4.350,0.660)--(4.350,0.785);
\draw[gp path] (4.350,2.706)--(4.350,2.581);
\gpcolor{color=gp lt color axes}
\gpsetlinetype{gp lt axes}
\draw[gp path] (4.551,0.660)--(4.551,2.706);
\gpcolor{color=gp lt color border}
\gpsetlinetype{gp lt border}
\draw[gp path] (4.551,0.660)--(4.551,0.785);
\draw[gp path] (4.551,2.706)--(4.551,2.581);
\gpcolor{color=gp lt color axes}
\gpsetlinetype{gp lt axes}
\draw[gp path] (4.714,0.660)--(4.714,2.706);
\gpcolor{color=gp lt color border}
\gpsetlinetype{gp lt border}
\draw[gp path] (4.714,0.660)--(4.714,0.785);
\draw[gp path] (4.714,2.706)--(4.714,2.581);
\gpcolor{color=gp lt color axes}
\gpsetlinetype{gp lt axes}
\draw[gp path] (4.853,0.660)--(4.853,2.706);
\gpcolor{color=gp lt color border}
\gpsetlinetype{gp lt border}
\draw[gp path] (4.853,0.660)--(4.853,0.785);
\draw[gp path] (4.853,2.706)--(4.853,2.581);
\gpcolor{color=gp lt color axes}
\gpsetlinetype{gp lt axes}
\draw[gp path] (4.972,0.660)--(4.972,2.706);
\gpcolor{color=gp lt color border}
\gpsetlinetype{gp lt border}
\draw[gp path] (4.972,0.660)--(4.972,0.785);
\draw[gp path] (4.972,2.706)--(4.972,2.581);
\gpcolor{color=gp lt color axes}
\gpsetlinetype{gp lt axes}
\draw[gp path] (5.078,0.660)--(5.078,2.706);
\gpcolor{color=gp lt color border}
\gpsetlinetype{gp lt border}
\draw[gp path] (5.078,0.660)--(5.078,0.785);
\draw[gp path] (5.078,2.706)--(5.078,2.581);
\gpcolor{color=gp lt color axes}
\gpsetlinetype{gp lt axes}
\draw[gp path] (5.173,0.660)--(5.173,2.706);
\gpcolor{color=gp lt color border}
\gpsetlinetype{gp lt border}
\draw[gp path] (5.173,0.660)--(5.173,0.911);
\draw[gp path] (5.173,2.706)--(5.173,2.455);
\gpcolor{rgb color={0.000,0.000,0.000}}
\node[gp node center,font={\fontsize{10pt}{12pt}\selectfont}] at (5.173,0.352) {$10^{0}$};
\gpcolor{color=gp lt color axes}
\gpsetlinetype{gp lt axes}
\draw[gp path] (5.795,0.660)--(5.795,2.218);
\draw[gp path] (5.795,2.526)--(5.795,2.706);
\gpcolor{color=gp lt color border}
\gpsetlinetype{gp lt border}
\draw[gp path] (5.795,0.660)--(5.795,0.785);
\draw[gp path] (5.795,2.706)--(5.795,2.581);
\gpcolor{color=gp lt color axes}
\gpsetlinetype{gp lt axes}
\draw[gp path] (6.159,0.660)--(6.159,2.218);
\draw[gp path] (6.159,2.526)--(6.159,2.706);
\gpcolor{color=gp lt color border}
\gpsetlinetype{gp lt border}
\draw[gp path] (6.159,0.660)--(6.159,0.785);
\draw[gp path] (6.159,2.706)--(6.159,2.581);
\gpcolor{color=gp lt color axes}
\gpsetlinetype{gp lt axes}
\draw[gp path] (6.417,0.660)--(6.417,2.218);
\draw[gp path] (6.417,2.526)--(6.417,2.706);
\gpcolor{color=gp lt color border}
\gpsetlinetype{gp lt border}
\draw[gp path] (6.417,0.660)--(6.417,0.785);
\draw[gp path] (6.417,2.706)--(6.417,2.581);
\gpcolor{color=gp lt color axes}
\gpsetlinetype{gp lt axes}
\draw[gp path] (6.617,0.660)--(6.617,2.218);
\draw[gp path] (6.617,2.526)--(6.617,2.706);
\gpcolor{color=gp lt color border}
\gpsetlinetype{gp lt border}
\draw[gp path] (6.617,0.660)--(6.617,0.785);
\draw[gp path] (6.617,2.706)--(6.617,2.581);
\gpcolor{color=gp lt color axes}
\gpsetlinetype{gp lt axes}
\draw[gp path] (6.781,0.660)--(6.781,2.218);
\draw[gp path] (6.781,2.526)--(6.781,2.706);
\gpcolor{color=gp lt color border}
\gpsetlinetype{gp lt border}
\draw[gp path] (6.781,0.660)--(6.781,0.785);
\draw[gp path] (6.781,2.706)--(6.781,2.581);
\gpcolor{color=gp lt color axes}
\gpsetlinetype{gp lt axes}
\draw[gp path] (6.919,0.660)--(6.919,2.218);
\draw[gp path] (6.919,2.526)--(6.919,2.706);
\gpcolor{color=gp lt color border}
\gpsetlinetype{gp lt border}
\draw[gp path] (6.919,0.660)--(6.919,0.785);
\draw[gp path] (6.919,2.706)--(6.919,2.581);
\gpcolor{color=gp lt color axes}
\gpsetlinetype{gp lt axes}
\draw[gp path] (7.039,0.660)--(7.039,2.218);
\draw[gp path] (7.039,2.526)--(7.039,2.706);
\gpcolor{color=gp lt color border}
\gpsetlinetype{gp lt border}
\draw[gp path] (7.039,0.660)--(7.039,0.785);
\draw[gp path] (7.039,2.706)--(7.039,2.581);
\gpcolor{color=gp lt color axes}
\gpsetlinetype{gp lt axes}
\draw[gp path] (7.144,0.660)--(7.144,2.706);
\gpcolor{color=gp lt color border}
\gpsetlinetype{gp lt border}
\draw[gp path] (7.144,0.660)--(7.144,0.785);
\draw[gp path] (7.144,2.706)--(7.144,2.581);
\gpcolor{color=gp lt color axes}
\gpsetlinetype{gp lt axes}
\draw[gp path] (7.239,0.660)--(7.239,2.706);
\gpcolor{color=gp lt color border}
\gpsetlinetype{gp lt border}
\draw[gp path] (7.239,0.660)--(7.239,0.911);
\draw[gp path] (7.239,2.706)--(7.239,2.455);
\gpcolor{rgb color={0.000,0.000,0.000}}
\node[gp node center,font={\fontsize{10pt}{12pt}\selectfont}] at (7.239,0.352) {$10^{1}$};
\gpcolor{color=gp lt color border}
\draw[gp path] (1.040,2.706)--(1.040,0.660)--(7.239,0.660)--(7.239,2.706)--cycle;
\gpcolor{rgb color={0.000,0.000,0.000}}
\node[gp node center,rotate=90,font={\fontsize{10pt}{12pt}\selectfont}] at (-0.218,1.683) {Phase [deg]};
\node[gp node center,font={\fontsize{10pt}{12pt}\selectfont}] at (4.139,-0.110) {Frequency [rad/s]};
\gpcolor{color=gp lt color border}
\gpsetlinewidth{1.00}
\draw[gp path] (5.403,2.218)--(5.403,2.526)--(7.055,2.526)--(7.055,2.218)--cycle;
\gpcolor{rgb color={0.000,0.000,0.000}}
\node[gp node right,font={\fontsize{10pt}{12pt}\selectfont}] at (7.055,2.372) {$G_1$};
\gpcolor{rgb color={0.000,0.000,1.000}}
\gpsetlinetype{gp lt plot 0}
\gpsetlinewidth{0.50}
\draw[gp path] (5.587,2.372)--(6.503,2.372);
\draw[gp path] (1.040,2.613)--(1.052,2.613)--(1.065,2.613)--(1.077,2.612)--(1.090,2.612)%
  --(1.102,2.612)--(1.115,2.612)--(1.127,2.612)--(1.139,2.612)--(1.152,2.611)--(1.164,2.611)%
  --(1.177,2.611)--(1.189,2.611)--(1.201,2.611)--(1.214,2.610)--(1.226,2.610)--(1.239,2.610)%
  --(1.251,2.610)--(1.264,2.610)--(1.276,2.609)--(1.288,2.609)--(1.301,2.609)--(1.313,2.609)%
  --(1.326,2.608)--(1.338,2.608)--(1.351,2.608)--(1.363,2.608)--(1.375,2.608)--(1.388,2.607)%
  --(1.400,2.607)--(1.413,2.607)--(1.425,2.607)--(1.438,2.606)--(1.450,2.606)--(1.462,2.606)%
  --(1.475,2.605)--(1.487,2.605)--(1.500,2.605)--(1.512,2.605)--(1.524,2.604)--(1.537,2.604)%
  --(1.549,2.604)--(1.562,2.603)--(1.574,2.603)--(1.587,2.603)--(1.599,2.603)--(1.611,2.602)%
  --(1.624,2.602)--(1.636,2.602)--(1.649,2.601)--(1.661,2.601)--(1.674,2.601)--(1.686,2.600)%
  --(1.698,2.600)--(1.711,2.600)--(1.723,2.599)--(1.736,2.599)--(1.748,2.598)--(1.761,2.598)%
  --(1.773,2.598)--(1.785,2.597)--(1.798,2.597)--(1.810,2.597)--(1.823,2.596)--(1.835,2.596)%
  --(1.847,2.595)--(1.860,2.595)--(1.872,2.595)--(1.885,2.594)--(1.897,2.594)--(1.910,2.593)%
  --(1.922,2.593)--(1.934,2.592)--(1.947,2.592)--(1.959,2.591)--(1.972,2.591)--(1.984,2.590)%
  --(1.997,2.590)--(2.009,2.589)--(2.021,2.589)--(2.034,2.588)--(2.046,2.588)--(2.059,2.587)%
  --(2.071,2.587)--(2.084,2.586)--(2.096,2.586)--(2.108,2.585)--(2.121,2.585)--(2.133,2.584)%
  --(2.146,2.583)--(2.158,2.583)--(2.170,2.582)--(2.183,2.582)--(2.195,2.581)--(2.208,2.580)%
  --(2.220,2.580)--(2.233,2.579)--(2.245,2.578)--(2.257,2.578)--(2.270,2.577)--(2.282,2.576)%
  --(2.295,2.576)--(2.307,2.575)--(2.320,2.574)--(2.332,2.574)--(2.344,2.573)--(2.357,2.572)%
  --(2.369,2.571)--(2.382,2.571)--(2.394,2.570)--(2.407,2.569)--(2.419,2.568)--(2.431,2.567)%
  --(2.444,2.567)--(2.456,2.566)--(2.469,2.565)--(2.481,2.564)--(2.493,2.563)--(2.506,2.562)%
  --(2.518,2.561)--(2.531,2.560)--(2.543,2.559)--(2.556,2.558)--(2.568,2.557)--(2.580,2.556)%
  --(2.593,2.555)--(2.605,2.554)--(2.618,2.553)--(2.630,2.552)--(2.643,2.551)--(2.655,2.550)%
  --(2.667,2.549)--(2.680,2.548)--(2.692,2.547)--(2.705,2.546)--(2.717,2.544)--(2.730,2.543)%
  --(2.742,2.542)--(2.754,2.541)--(2.767,2.539)--(2.779,2.538)--(2.792,2.537)--(2.804,2.535)%
  --(2.816,2.534)--(2.829,2.533)--(2.841,2.531)--(2.854,2.530)--(2.866,2.528)--(2.879,2.527)%
  --(2.891,2.525)--(2.903,2.524)--(2.916,2.522)--(2.928,2.521)--(2.941,2.519)--(2.953,2.517)%
  --(2.966,2.516)--(2.978,2.514)--(2.990,2.512)--(3.003,2.510)--(3.015,2.508)--(3.028,2.507)%
  --(3.040,2.505)--(3.053,2.503)--(3.065,2.501)--(3.077,2.499)--(3.090,2.497)--(3.102,2.495)%
  --(3.115,2.492)--(3.127,2.490)--(3.139,2.488)--(3.152,2.486)--(3.164,2.483)--(3.177,2.481)%
  --(3.189,2.478)--(3.202,2.476)--(3.214,2.473)--(3.226,2.471)--(3.239,2.468)--(3.251,2.465)%
  --(3.264,2.463)--(3.276,2.460)--(3.289,2.457)--(3.301,2.454)--(3.313,2.451)--(3.326,2.447)%
  --(3.338,2.444)--(3.351,2.441)--(3.363,2.437)--(3.375,2.434)--(3.388,2.430)--(3.400,2.426)%
  --(3.413,2.423)--(3.425,2.419)--(3.438,2.415)--(3.450,2.411)--(3.462,2.406)--(3.475,2.402)%
  --(3.487,2.397)--(3.500,2.393)--(3.512,2.388)--(3.525,2.383)--(3.537,2.378)--(3.549,2.372)%
  --(3.562,2.367)--(3.574,2.361)--(3.587,2.356)--(3.599,2.350)--(3.612,2.343)--(3.624,2.337)%
  --(3.636,2.330)--(3.649,2.323)--(3.661,2.316)--(3.674,2.309)--(3.686,2.301)--(3.698,2.293)%
  --(3.711,2.285)--(3.723,2.276)--(3.736,2.267)--(3.748,2.258)--(3.761,2.248)--(3.773,2.238)%
  --(3.785,2.227)--(3.798,2.216)--(3.810,2.205)--(3.823,2.193)--(3.835,2.181)--(3.848,2.168)%
  --(3.860,2.154)--(3.872,2.140)--(3.885,2.126)--(3.897,2.110)--(3.910,2.094)--(3.922,2.078)%
  --(3.935,2.061)--(3.947,2.043)--(3.959,2.024)--(3.972,2.005)--(3.984,1.985)--(3.997,1.964)%
  --(4.009,1.942)--(4.021,1.920)--(4.034,1.897)--(4.046,1.874)--(4.059,1.849)--(4.071,1.825)%
  --(4.084,1.800)--(4.096,1.774)--(4.108,1.748)--(4.121,1.722)--(4.133,1.696)--(4.139,1.683)%
  --(4.146,1.670)--(4.158,1.644)--(4.171,1.618)--(4.183,1.592)--(4.195,1.566)--(4.208,1.541)%
  --(4.220,1.517)--(4.233,1.492)--(4.245,1.469)--(4.258,1.446)--(4.270,1.424)--(4.282,1.402)%
  --(4.295,1.381)--(4.307,1.361)--(4.320,1.342)--(4.332,1.323)--(4.344,1.305)--(4.357,1.288)%
  --(4.369,1.272)--(4.382,1.256)--(4.394,1.240)--(4.407,1.226)--(4.419,1.212)--(4.431,1.198)%
  --(4.444,1.185)--(4.456,1.173)--(4.469,1.161)--(4.481,1.150)--(4.494,1.139)--(4.506,1.128)%
  --(4.518,1.118)--(4.531,1.108)--(4.543,1.099)--(4.556,1.090)--(4.568,1.081)--(4.581,1.073)%
  --(4.593,1.065)--(4.605,1.057)--(4.618,1.050)--(4.630,1.043)--(4.643,1.036)--(4.655,1.029)%
  --(4.667,1.023)--(4.680,1.016)--(4.692,1.010)--(4.705,1.005)--(4.717,0.999)--(4.730,0.994)%
  --(4.742,0.988)--(4.754,0.983)--(4.767,0.978)--(4.779,0.973)--(4.792,0.969)--(4.804,0.964)%
  --(4.817,0.960)--(4.829,0.955)--(4.841,0.951)--(4.854,0.947)--(4.866,0.943)--(4.879,0.940)%
  --(4.891,0.936)--(4.904,0.932)--(4.916,0.929)--(4.928,0.925)--(4.941,0.922)--(4.953,0.919)%
  --(4.966,0.915)--(4.978,0.912)--(4.990,0.909)--(5.003,0.906)--(5.015,0.903)--(5.028,0.901)%
  --(5.040,0.898)--(5.053,0.895)--(5.065,0.893)--(5.077,0.890)--(5.090,0.888)--(5.102,0.885)%
  --(5.115,0.883)--(5.127,0.880)--(5.140,0.878)--(5.152,0.876)--(5.164,0.874)--(5.177,0.871)%
  --(5.189,0.869)--(5.202,0.867)--(5.214,0.865)--(5.226,0.863)--(5.239,0.861)--(5.251,0.859)%
  --(5.264,0.858)--(5.276,0.856)--(5.289,0.854)--(5.301,0.852)--(5.313,0.850)--(5.326,0.849)%
  --(5.338,0.847)--(5.351,0.845)--(5.363,0.844)--(5.376,0.842)--(5.388,0.841)--(5.400,0.839)%
  --(5.413,0.838)--(5.425,0.836)--(5.438,0.835)--(5.450,0.833)--(5.463,0.832)--(5.475,0.831)%
  --(5.487,0.829)--(5.500,0.828)--(5.512,0.827)--(5.525,0.825)--(5.537,0.824)--(5.549,0.823)%
  --(5.562,0.822)--(5.574,0.820)--(5.587,0.819)--(5.599,0.818)--(5.612,0.817)--(5.624,0.816)%
  --(5.636,0.815)--(5.649,0.814)--(5.661,0.813)--(5.674,0.812)--(5.686,0.811)--(5.699,0.810)%
  --(5.711,0.809)--(5.723,0.808)--(5.736,0.807)--(5.748,0.806)--(5.761,0.805)--(5.773,0.804)%
  --(5.786,0.803)--(5.798,0.802)--(5.810,0.801)--(5.823,0.800)--(5.835,0.799)--(5.848,0.799)%
  --(5.860,0.798)--(5.872,0.797)--(5.885,0.796)--(5.897,0.795)--(5.910,0.795)--(5.922,0.794)%
  --(5.935,0.793)--(5.947,0.792)--(5.959,0.792)--(5.972,0.791)--(5.984,0.790)--(5.997,0.790)%
  --(6.009,0.789)--(6.022,0.788)--(6.034,0.788)--(6.046,0.787)--(6.059,0.786)--(6.071,0.786)%
  --(6.084,0.785)--(6.096,0.784)--(6.109,0.784)--(6.121,0.783)--(6.133,0.783)--(6.146,0.782)%
  --(6.158,0.781)--(6.171,0.781)--(6.183,0.780)--(6.195,0.780)--(6.208,0.779)--(6.220,0.779)%
  --(6.233,0.778)--(6.245,0.778)--(6.258,0.777)--(6.270,0.777)--(6.282,0.776)--(6.295,0.776)%
  --(6.307,0.775)--(6.320,0.775)--(6.332,0.774)--(6.345,0.774)--(6.357,0.773)--(6.369,0.773)%
  --(6.382,0.772)--(6.394,0.772)--(6.407,0.771)--(6.419,0.771)--(6.432,0.771)--(6.444,0.770)%
  --(6.456,0.770)--(6.469,0.769)--(6.481,0.769)--(6.494,0.769)--(6.506,0.768)--(6.518,0.768)%
  --(6.531,0.768)--(6.543,0.767)--(6.556,0.767)--(6.568,0.766)--(6.581,0.766)--(6.593,0.766)%
  --(6.605,0.765)--(6.618,0.765)--(6.630,0.765)--(6.643,0.764)--(6.655,0.764)--(6.668,0.764)%
  --(6.680,0.763)--(6.692,0.763)--(6.705,0.763)--(6.717,0.763)--(6.730,0.762)--(6.742,0.762)%
  --(6.755,0.762)--(6.767,0.761)--(6.779,0.761)--(6.792,0.761)--(6.804,0.761)--(6.817,0.760)%
  --(6.829,0.760)--(6.841,0.760)--(6.854,0.759)--(6.866,0.759)--(6.879,0.759)--(6.891,0.759)%
  --(6.904,0.758)--(6.916,0.758)--(6.928,0.758)--(6.941,0.758)--(6.953,0.758)--(6.966,0.757)%
  --(6.978,0.757)--(6.991,0.757)--(7.003,0.757)--(7.015,0.756)--(7.028,0.756)--(7.040,0.756)%
  --(7.053,0.756)--(7.065,0.756)--(7.078,0.755)--(7.090,0.755)--(7.102,0.755)--(7.115,0.755)%
  --(7.127,0.755)--(7.140,0.754)--(7.152,0.754)--(7.164,0.754)--(7.177,0.754)--(7.189,0.754)%
  --(7.202,0.754)--(7.214,0.753)--(7.227,0.753)--(7.239,0.753);
%% coordinates of the plot area
\gpdefrectangularnode{gp plot 2}{\pgfpoint{1.040cm}{0.660cm}}{\pgfpoint{7.239cm}{2.706cm}}
\end{tikzpicture}
%% gnuplot variables

      \centering
      \caption{Diagrama de Bode para a função de transferência $G_1(s)=\frac{10}{s^2+2s+1}$}
    \end{figure}
\end{enumerate}
\end{document}
